\section{Parametric systems and classification}

Some systems contain a parameter. Then the question is not only ``What is the solution?'' but also ``For which parameter values does the structure of the system change?''

This is where determinants, row reduction, and geometric interpretation come together.

\subsection*{The three possible outcomes}

Every linear system has one of the following outcomes:
\begin{itemize}
\item exactly one solution,
\item no solution,
\item infinitely many solutions.
\end{itemize}

For a square system
\[
Ax=b,
\]
one important test is the determinant of \(A\). If
\[
\det A\ne0,
\]
then \(A\) is invertible and the system has exactly one solution. If
\[
\det A=0,
\]
then the system does not have a unique solution obtained by inversion. But we must still decide whether it has no solution or infinitely many solutions.

\begin{remarkbox}
The condition \(\det A=0\) is a warning sign, not a complete answer. It says that uniqueness is lost. It does not by itself decide whether solutions exist.
\end{remarkbox}

\subsection*{A simple parameter example}

Consider the system
\[
\begin{cases}
x+y=1,\\
2x+2y=a.
\end{cases}
\]
The left-hand side of the second equation is twice the left-hand side of the first equation. Therefore consistency depends on the right-hand side.

If we multiply the first equation by \(2\), we get
\[
2x+2y=2.
\]
So for the second equation to describe the same condition, we need
\[
a=2.
\]

If \(a=2\), the two equations describe the same line, and there are infinitely many solutions:
\[
x+y=1.
\]
If \(a\ne2\), the equations describe parallel distinct lines, and there is no solution.

\begin{summarybox}
For
\[
\begin{cases}
x+y=1,\\
2x+2y=a,
\end{cases}
\]
we have:
\begin{itemize}
\item if \(a=2\), infinitely many solutions,
\item if \(a\ne2\), no solution.
\end{itemize}
\end{summarybox}

This example shows that a parameter may affect only the right-hand side, but that can still change the solution set completely.

\subsection*{Using determinants to find exceptional values}

Now consider a parameter in the coefficient matrix. Let
\[
A(t)=\begin{pmatrix}
t&1\\
4&t
\end{pmatrix}.
\]
The determinant is
\[
\det A(t)=t^2-4.
\]
The exceptional values occur when
\[
t^2-4=0,
\]
so
\[
t=2
\qquad\text{or}\qquad
 t=-2.
\]
For all other values of \(t\), the matrix is invertible. Therefore any system
\[
A(t)x=b
\]
has exactly one solution for \(t\ne2\) and \(t\ne-2\).

The exceptional values \(t=2\) and \(t=-2\) require separate analysis. At those values, the determinant test tells us that the system is not uniquely solvable, but not whether it is inconsistent or has infinitely many solutions.

\subsection*{A full classification example}

\begin{examplebox}
Classify the system depending on the parameter \(a\):
\[
\begin{cases}
x+y=3,\\
2x+ay=6.
\end{cases}
\]
The coefficient matrix is
\[
A(a)=\begin{pmatrix}
1&1\\
2&a
\end{pmatrix}.
\]
Its determinant is
\[
\det A(a)=a-2.
\]
If
\[
a\ne2,
\]
then \(\det A(a)\ne0\), so the system has exactly one solution.

Now consider the exceptional value \(a=2\). The system becomes
\[
\begin{cases}
x+y=3,\\
2x+2y=6.
\end{cases}
\]
The second equation is twice the first. Therefore the system has infinitely many solutions.

Thus:
\begin{itemize}
\item if \(a\ne2\), the system has exactly one solution,
\item if \(a=2\), the system has infinitely many solutions.
\end{itemize}
\end{examplebox}

The determinant found the value where the structure changes. Row interpretation then classified what happens at that value.

\begin{examplebox}
Classify the system depending on \(a\):
\[
\begin{cases}
x+y=3,\\
2x+2y=a.
\end{cases}
\]
The coefficient matrix is
\[
A=\begin{pmatrix}
1&1\\
2&2
\end{pmatrix},
\]
so
\[
\det A=0.
\]
There is no value of \(a\) that makes this coefficient matrix invertible. We must compare the equations.

Twice the first equation gives
\[
2x+2y=6.
\]
Therefore:
\begin{itemize}
\item if \(a=6\), the two equations describe the same line, so there are infinitely many solutions,
\item if \(a\ne6\), the two equations are inconsistent, so there is no solution.
\end{itemize}
\end{examplebox}

This example shows that a singular coefficient matrix may still lead to different outcomes depending on the right-hand side.

\subsection*{Parameters and row reduction}

For larger systems, row reduction is often the safest method. It allows us to see contradictions, zero rows, and free variables directly.

A row such as
\[
0=0
\]
indicates redundancy. A row such as
\[
0=5
\]
indicates inconsistency. A missing pivot indicates a free variable. These patterns remain valid even when parameters are present.

The difficulty is that some row operations may require dividing by an expression containing a parameter. Before dividing by such an expression, we must consider when it might be zero.

\begin{remarkbox}
When a parameter appears, do not divide by an expression such as \(a-2\) without asking whether \(a-2=0\) is possible. Parameter values that make a denominator zero are usually exceptional cases and must be treated separately.
\end{remarkbox}

\subsection*{A classification strategy}

A good solution to a parameter system should not only list parameter values. It should explain why the classification changes.

\begin{summarybox}
For a square parameter system \(A(t)x=b(t)\), a useful strategy is:
\begin{enumerate}
\item Compute \(\det A(t)\).
\item Solve \(\det A(t)=0\) to find exceptional parameter values.
\item For non-exceptional values, conclude that there is exactly one solution.
\item For exceptional values, analyze the system separately using row reduction or equation comparison.
\item Classify each case as exactly one solution, no solution, or infinitely many solutions.
\end{enumerate}
\end{summarybox}

The determinant identifies where uniqueness can fail. Elimination and interpretation decide what happens there.

\subsection*{What the chapter has built}

Systems of linear equations bring together the earlier algebraic tools.

Matrices organize the coefficients. Row operations simplify the system. Determinants detect whether a square coefficient matrix is invertible. Inverse matrices give the formula \(x=A^{-1}b\) when the inverse exists. Cramer's rule gives a determinant-based expression for the solution when \(\det A\ne0\).

These methods are not isolated tricks. They are different ways of reading the same object: a system of linear conditions.

\begin{summarybox}
At the end of this chapter, remember:
\begin{itemize}
\item A system is a set of conditions to be satisfied simultaneously.
\item The augmented matrix records the coefficients and the right-hand side.
\item Gaussian elimination is the most flexible method for solving and classifying systems.
\item If \(\det A\ne0\), a square system \(Ax=b\) has exactly one solution.
\item If \(\det A=0\), uniqueness is lost, but existence must still be checked.
\item A linear system may have exactly one solution, no solution, or infinitely many solutions.
\end{itemize}
\end{summarybox}

The goal is not only to solve systems. The goal is to understand what the equations say about structure, dependence, and compatibility.